\documentclass{article}
\usepackage{graphicx} % Required for inserting images
\usepackage[a4paper, left=2cm, right=2cm, top=3cm, bottom=3cm]{geometry}
\usepackage{amsmath}
\usepackage{amssymb}
\usepackage{centernot}
\allowdisplaybreaks

\title{ICLR Electronics Notes}
\author{Siheon Lee}
\date{December 2025}

\begin{document}

\maketitle

\section{Sensors (06/12/2025)}
\begin{itemize}
    \item Sensor fusion to estimate orientation/attitude/heading
    \item Fusion algorithm can be referred to as attitude and heading reference system (AHRS)
    \item To define an orientation, we need a reference frame and specify rotation
    \item Represent 3D rotation between two coordinate frames
    \begin{itemize}
        \item Roll, pitch, yaw
        \item Direction cosine matrix (DCM)
        \item Quaternion
    \end{itemize}
    \item Magnetometer
    \begin{itemize}
        \item Points towards north but is 3D so returns 3D vector in direction of gravity
        \item Affected by disturbances in the magnetic field
        \begin{itemize}
            \item If magnetic disturbance is part of system and rotates with magnetometer, it can be calibrated out
            \item Hard iron source - generates own magnetic field (magnet, coil, etc.)
            \begin{itemize}
                \item If we rotate the system around a single axis and measure the magnetic field, we will find circle offset from origin
                \item Magnetometer read larger intensity in one direction and smaller intensity in opposite direction
            \end{itemize}
            \item Soft iron source - something that doesn't generate its own magnetic field but bends magnetic field (iron nail, structural elements, etc.)
            \begin{itemize}
                \item Distort measurement as an oval
            \end{itemize}
        \end{itemize}
        \item Magnetometer calibration
        \begin{itemize}
            \item Measuring across $4 \pi$ steradian directions, we would find perfect sphere in perfect system with radius being magnitude of sphere
            \item By measuring the distortions, we can find a correction by $x_{corrected} = (x - b) \cdot A$ where $b$ is the hard iron bias ($3 \times 1$ vector) and $A$ is the soft iron distortion ($3 \times 3$ matrix)
        \end{itemize}
    \end{itemize}
    \item Accelerometer
    \begin{itemize}
        \item Down is oppoosite direction of acceleration vector 
        \item East is cross product of down and measured magnetic field
        \item North is cross-product of east and down
        \item Direction cosine matrix can be built from north, east, and down vectors
        \item Accelerometer not centered at center of rotation will sense acceleration
        \begin{itemize}
            \item Predicting linear acceleration and removing it from measurement prior to using it
            \item If acceleration is result of actuators, we can run it through a system model to get an estimated linear acceleration and subtract it from the accelerometer measurement
            \item We can also ignore $\vert a - g \vert > t$ for some acceleration measurement $a$, gravity $g$, and threshold $t$
        \end{itemize}
    \end{itemize}
    \item Gyroscope
    \begin{itemize}
        \item $a \cdot \Delta t = \Delta a$ where $a$ is measured angular rate, $\Delta t$ is sample time, and $\Delta a$ is change in angle
        \item Add $\Delta a$ to previous orientation to get current orientation
        \begin{itemize}
            \item Dead reckoning and just is integrating gyro measurement
            \item Downsides: need to known initial orientation and sensor biases
            \item Integration acts as a low pass filter so that high frequency noise is smoothed out but result drifts away from true position due to random walk and integrating any biases
        \end{itemize}
    \end{itemize}
    \item Inertial measurement unit often packages gyroscope and accelerometer
    \item Sensor fusion
    \begin{itemize}
        \item Accel + mag
        \begin{itemize}
            \item Produces absolute measurement
            \item Corrupted by common disturbances
        \end{itemize}
        \item Gyro
        \begin{itemize}
            \item Produces relative measurement
            \item Needs initial orientation
            \item Drifts over time
        \end{itemize}
        \item Sensor fusion algorithsm: complementary, Kalman, Madgwick, Mahony filters
        \begin{itemize}
            \item At basic, initializes attitude (manual or initial results of mag and accelerometer) and uses direction of mag field and gravity to slowly correct for drift in gyro
        \end{itemize}
    \end{itemize}
    \item The degree to which to use accel+mag and gyro measurements: complementary filter is decided by yourself while for the Kalman filter, it is calculated for you (specify noise and how good system model is)
    \item Accelerometer, gyroscope, barometer
\end{itemize}

\section{(09/12/2025)}
\begin{itemize}
    \item https://jlcpcb.com/parts/
    \item https://componentsearchengine.com/
    \item Download model, unzip, find kicad folder, and put it in working directory
\end{itemize}

\section{DPS 368 Modeling (17/01/2026)}
\subsection{Workflow}
\begin{itemize}
    \item Read from 0x10 to 0x20 (COEF)
    \item Write to 0x06 (PRS\_CFG)
    \item Write to 0x07 (TMP\_CFG)
    \item Write to 0x09 (CFG\_REG)
    \item Loop?
    \begin{itemize}
        \item Write 0x02 to 0x08 (PRS\_B0, TMP\_B2, TMP\_B1, TMP\_B0, PRS\_CFG, TMP\_CFG, MEAS\_CFG)
        \item Read 0x03 to 0x05 (TMP\_B2, TMP\_B1, TMP\_B0)
        \item Calculate Temperature
        \item Write 0x01 to 0x08 (PRS\_B1, PRS\_B0, TMP\_B2, TMP\_B1, TMP\_B0, PRS\_CFG, TMP\_CFG, MEAS\_CFG)
        \item Read 0x00 to 0x02 (PRS\_B2, PRS\_B1, PRS\_B0)
        \item Calculate Pressure
    \end{itemize}
    \item End
\end{itemize}
\subsection{General Notes}
\begin{itemize}
    \item Oversampling
    \begin{itemize}
        \item Takes multiple raw pressure measurements
        \item Averages (or digitally filters) them
        \item Outputs a single pressure value with lower noise and higher effective resolution
        \item Set oversampling rate through TMP\_PRC[2:0] (temperature) and PRS\_PRC[3:0] (pressure)
    \end{itemize}
    \item FIFO
    \begin{itemize}
        \item Stores up to 32 pressure or temperature measurements.
        \item Measurement time depends on oversampling rate (27.6 ms for standard mode 16x)
        \item Read the pressure calibration coefficients (c00, c10, c20, c30, c01, c11, and c21) from the Calibration Coefficient register
        \item Choose scaling factors kT (for temperature) and kP (for pressure) based on the chosen precision rate
        \item Read the pressure and temperature result from the registers or FIFO
        \item $T_{\text{raw\_sc}}=\frac{T_{\text{raw}}}{kT}$
    \end{itemize}
    \item Compensated scale factors are calculated based on the oversampling rate
    \item How to calculate compensated temperature values
    \begin{itemize}
        \item Read the temperature calibration coefficients ( c0 and c1 ) from the Calibration Coefficients (COEF) register
        \item Choose scaling factors kT (for temperature) based on the chosen precision rate
        \item Read the temperature result from the temperature register or FIFO
        \item Calculate scaled measurement results: $T_{\text{raw\_sc}}=\frac{T_{\text{raw}}}{kT}$
        \item Calculate compensated measurement results: $T_{\text{comp}} = (c0)0.5 + (c1)(T_{\text{raw\_sc}})$
    \end{itemize}
    \item How to calculate compensated pressure values
    \begin{itemize}
        \item Read the pressure calibration coefficients (c00, c10, c20, c30, c01, c11, and c21) from the Calibration Coefficient register
        \item Choose scaling factors kT (for temperature) and kP (for pressure) based on the chosen precision rate
        \item Read the pressure and temperature result from the registers or FIFO
        \item Calculate scaled measurement results:
        \begin{flalign*}
        T_{\text{raw\_sc}}=\frac{T_{\text{raw}}}{kT} \\
        T_{\text{raw\_sc}}=\frac{T_{\text{raw}}}{kT} \\
        \end{flalign*}
        \item Calculate compensated measurement results: $T_{\text{comp}} = (c0)0.5 + (c1)(T_{\text{raw\_sc}})$
    \end{itemize}
    \item Barometric pressure (with temperature gradient)
    \begin{itemize}
        \item $P(h)$: pressure at 
        \item $P_b$: reference pressure
        \item $T_{M,b}$: reference temperature
        \item $L_{M,b}$: temperature gradient
        \item $H$: current height
        \item $H_b$: reference height
        \item $R^*$: universal gas constant
        \item $M_0$: mean molar mass of air at sea level
        \item $g_0'$: gravitational acceleration
        \begin{align*}
            P(H) &= P_b \cdot \left[\frac{T_{M,b}}{T_{M,b} + L_{M,b} \cdot (H-H_b)}\right]^{\frac{g_0' \cdot M_0}{R^* \cdot L_{M,b}}} \\
            H &= \frac{T_{M,b}}{L_{M,b}} \left[ \left(\frac{P(H)}{P_b}\right)^{-\frac{R^* \cdot L_{M,b}}{g_0' \cdot M_0}} - 1 \right] + H_b
        \end{align*}
    \end{itemize}
    \item Discrete-Time Gauss-Markov Drift Model
    \begin{itemize}
        \item $\omega_k \sim \mathcal{N}(0,q)$: process noise
        \item $\tau$: correlation time
        \begin{align*}
            b_{k+1} &= \phi b_k + \omega_k \\
            \phi &= e^{-\Delta t/ \tau}
        \end{align*}
    \end{itemize}
\end{itemize}

\end{document}

